% !TEX TS-program = lualatex
Диаграмма \refimage{model-builder-usecase} иллюстрирует варианты использования модуля для тренировки интеллектуальных агентов --- инструмента для построения и обучения интеллектуальных агентов на основе модели TD-Gammon. Взаимодействие с системой осуществляется двумя основными субъектами: пользователем, инициирующим процесс обучения, и системой, которая автоматически выполняет определённые действия, связанные с обучением и сохранением модели.

\image[width=0.9\textwidth]{Схема использования модуля для тренировки интеллектуальных агентов}{model-builder-usecase}{model-builder-usecase}

Процесс начинается с запуска пользователем модуля построителя моделей. После запуска пользователь указывает параметры обучения, включая скорость обучения (\texttt{-{}-}lr), количество нейронов в скрытом слое (\texttt{-{}-}hidden\_units), параметр $\lambda$ для следов посещения (\texttt{-{}-}lamda), количество эпизодов обучения (\texttt{-{}-}episodes), начальное значение генератора случайных чисел (\texttt{-{}-}seed) и другие. Это действие представлено как вариант использования.

Далее пользователь может либо создать новую модель TD-Gammon, либо загрузить ранее сохранённую. Возможность загрузки модели является расширением основного сценария, применимым в случае указания аргумента \texttt{-{}-}model. При создании новой модели выполняются два шага: инициализация весов и настройка окружения Backgammon. Если же модель загружается, происходит только настройка окружения.

После подготовки среды и модели начинается этап обучения. Этот вариант использования выполняется как по инициативе пользователя, так и со стороны самой системы. При наличии указания \texttt{-{}-}save\_path в параметрах запуска, процесс включает сохранение результатов --- веса модели и параметры эксперимента сохраняются для последующего использования или анализа. Этот шаг также может инициироваться как пользователем, так и системой.
